% Source: Wald GR, physical PDF pp. 452--454 (printed pp. 439--441).
% Coverage: Appendix C.2, Lie Derivatives; equations (C.2.1)--(C.2.17).
% Figures/tables/footnotes: none.
% Status: source-reviewed and compile-clean.
% Uncertainties: none.

\subsection{Lie Derivatives}\label{sec:C.2}

Let \(M\) be a manifold and let \(\phi_t\) be a one-parameter group of
diffeomorphisms. As discussed in Section~2.2, \(\phi_t\) will be generated by a
vector field, \(v^a\). By the results of the previous section, we can use
\(\phi_t^*\) to carry along a smooth tensor field
\(T^{a_1\ldots a_k}{}_{b_1\ldots b_l}\). Comparison of
\(T^{a_1\ldots a_k}{}_{b_1\ldots b_l}\) and
\(\phi_{-t}^*T^{a_1\ldots a_k}{}_{b_1\ldots b_l}\) for small \(t\) gives rise to
the notion of the Lie derivative, \(\Lie_v\), with respect to \(v^a\). More
precisely, we define \(\Lie_v\) by
\begin{equation}
  \Lie_vT^{a_1\ldots a_k}{}_{b_1\ldots b_l}
  =\lim_{t\to0}\left\{
    \frac{(\phi_{-t}^*T)^{a_1\ldots a_k}{}_{b_1\ldots b_l}
    -T^{a_1\ldots a_k}{}_{b_1\ldots b_l}}{t}
  \right\},
  \tag{C.2.1}\label{eq:C.2.1}
\end{equation}
where all tensors appearing in Eq.~\eqref{eq:C.2.1} are evaluated at the same point
\(p\). Note that the vector index on \(v^a\) is dropped in the symbol \(\Lie_v\)
since its presence could lead to confusion.

It follows immediately from its definition, Eq.~\eqref{eq:C.2.1}, that \(\Lie_v\)
is a linear map from smooth tensor fields of type \((k,l)\) to smooth tensor fields of
type \((k,l)\). It also is not difficult to show (see Eq.~\eqref{eq:C.2.4} below)
that \(\Lie_v\) satisfies the Leibniz rule on outer products of tensors. Furthermore,
since \(v^a\) is tangent to the integral curves of \(\phi_t\), for functions
\(f:M\to\mathbb R\) we have
\begin{equation}
  \Lie_v(f)=v(f).
  \tag{C.2.2}\label{eq:C.2.2}
\end{equation}
Note also that
\(\Lie_vT^{a_1\ldots a_k}{}_{b_1\ldots b_l}=0\) everywhere if and only if for all
\(t\), \(\phi_t\) is a symmetry transformation for
\(T^{a_1\ldots a_k}{}_{b_1\ldots b_l}\).

To analyze the action of \(\Lie_v\) on an arbitrary tensor field, it is helpful to
introduce a coordinate system on \(M\) where the parameter \(t\) along the integral
curves of \(v^a\) is chosen as one of the coordinates \(x^1\), so that
\(v^a=(\partial/\partial x^1)^a\). (This always can be done locally in any region
where \(v^a\ne0\).) The action of \(\phi_{-t}\) then corresponds to the coordinate
transformation \(x^1\to x^1+t\), with \(x^2,\ldots,x^n\) held fixed. From the
parenthetical remark below Eq.~\eqref{eq:C.1.1}, we have
\((\phi^*)^\mu{}_\nu=\delta^\mu{}_\nu\) and hence, the coordinate-basis components
of \(\phi_{-t}^*T^{a_1\ldots a_k}{}_{b_1\ldots b_l}\) at the point \(p\) whose
coordinates are \((x^1,\ldots,x^n)\) are
\begin{equation}
  (\phi_{-t}^*T)^{\mu_1\ldots\mu_k}{}_{\nu_1\ldots\nu_l}
  (x^1,\ldots,x^n)
  =T^{\mu_1\ldots\mu_k}{}_{\nu_1\ldots\nu_l}
  (x^1+t,x^2,\ldots,x^n).
  \tag{C.2.3}\label{eq:C.2.3}
\end{equation}
Consequently, the components of the Lie derivative of
\(T^{a_1\ldots a_k}{}_{b_1\ldots b_l}\) in a coordinate system adapted to \(v^a\)
are simply
\begin{equation}
  \Lie_vT^{\mu_1\ldots\mu_k}{}_{\nu_1\ldots\nu_l}
  =\frac{\partial T^{\mu_1\ldots\mu_k}{}_{\nu_1\ldots\nu_l}}{\partial x^1}.
  \tag{C.2.4}\label{eq:C.2.4}
\end{equation}
Thus, in particular, \(\phi_t\) will be a symmetry transformation of
\(T^{a_1\ldots a_k}{}_{b_1\ldots b_l}\) if and only if the components
\(T^{\mu_1\ldots\mu_k}{}_{\nu_1\ldots\nu_l}\) in a coordinate system adapted to
\(v^a\) are independent of the integral curve coordinate \(x^1\).

We can obtain a coordinate-independent expression for the Lie derivative of a vector
field \(w^a\) by noting that in an adapted coordinate system we have by
Eq.~\eqref{eq:C.2.4},
\begin{equation}
  \Lie_vw^\mu=\frac{\partial w^\mu}{\partial x^1}.
  \tag{C.2.5}\label{eq:C.2.5}
\end{equation}
On the other hand, since \(v^a=(\partial/\partial x^1)^a\) and
\(w^a=\sum_\mu w^\mu(\partial/\partial x^\mu)^a\), the commutator of \(v^a\) and
\(w^a\) is given by
\begin{equation}
\begin{aligned}
  [v,w]^\mu
  &=\sum_\nu\left(
    v^\nu\frac{\partial w^\mu}{\partial x^\nu}
    -w^\nu\frac{\partial v^\mu}{\partial x^\nu}
  \right) \\
  &=\frac{\partial w^\mu}{\partial x^1}.
\end{aligned}
  \tag{C.2.6}\label{eq:C.2.6}
\end{equation}
Thus, we find that the components of \(\Lie_vw^a\) and \([v,w]^a\) are equal in an
adapted coordinate system. However, since both of these quantities are defined in a
coordinate-independent manner, we obtain
\begin{equation}
  \Lie_vw^a=[v,w]^a,
  \tag{C.2.7}\label{eq:C.2.7}
\end{equation}
which is the coordinate-independent formula we sought for the Lie derivative of a
vector field.

The action of the Lie derivative on all other types of tensor fields is determined by
Eqs.~\eqref{eq:C.2.2}, \eqref{eq:C.2.7} and the Leibniz rule. For example, for a
covector field, \(\mu_a\), we have by Eq.~\eqref{eq:C.2.2}
\begin{equation}
  \Lie_v(\mu_aw^a)=v(\mu_aw^a),
  \tag{C.2.8}\label{eq:C.2.8}
\end{equation}
where \(w^a\) is an arbitrary field. On the other hand, by the Leibniz rule and
Eq.~\eqref{eq:C.2.7}, we have
\begin{equation}
  \Lie_v(\mu_aw^a)=w^a\Lie_v\mu_a+\mu_a[v,w]^a.
  \tag{C.2.9}\label{eq:C.2.9}
\end{equation}
From the equality of the right sides of Eqs.~\eqref{eq:C.2.8} and \eqref{eq:C.2.9}
we obtain a formula which determines \(\Lie_v\mu_a\). This formula is most
conveniently expressed in terms of a derivative operator. If \(\nabla_a\) is an
arbitrary derivative operator on \(M\), we have by properties (4) and (2) of the
definition of derivative operator (see Section~3.1)
\begin{equation}
\begin{aligned}
  v(\mu_aw^a)
  &=v^b\nabla_b(\mu_aw^a)\\
  &=v^bw^a\nabla_b\mu_a+v^b\mu_a\nabla_bw^a.
\end{aligned}
  \tag{C.2.10}\label{eq:C.2.10}
\end{equation}
On the other hand, we showed previously (see Eq.~\eqref{eq:3.1.2}) that
\begin{equation}
  [v,w]^a=v^b\nabla_bw^a-w^b\nabla_bv^a.
  \tag{C.2.11}\label{eq:C.2.11}
\end{equation}
Thus, we find
\begin{equation}
  v^bw^a\nabla_b\mu_a+v^b\mu_a\nabla_bw^a
  =w^a\Lie_v\mu_a+\mu_av^b\nabla_bw^a-\mu_aw^b\nabla_bv^a,
  \tag{C.2.12}\label{eq:C.2.12}
\end{equation}
i.e.,
\begin{equation}
  \Lie_v\mu_a=v^b\nabla_b\mu_a+\mu_b\nabla_av^b.
  \tag{C.2.13}\label{eq:C.2.13}
\end{equation}
More generally for an arbitrary tensor field
\(T^{a_1\ldots a_k}{}_{b_1\ldots b_l}\) we find by induction that
\begin{equation}
\begin{aligned}
  \Lie_vT^{a_1\ldots a_k}{}_{b_1\ldots b_l}
  & =v^c\nabla_cT^{a_1\ldots a_k}{}_{b_1\ldots b_l}
  -\sum_{i=1}^{k}
  T^{a_1\ldots c\ldots a_k}{}_{b_1\ldots b_l}\nabla_cv^{a_i}\\
  &\quad +\sum_{j=1}^{l}
  T^{a_1\ldots a_k}{}_{b_1\ldots c\ldots b_l}\nabla_{b_j}v^c.
\end{aligned}
  \tag{C.2.14}\label{eq:C.2.14}
\end{equation}
Again, we emphasize that Eq.~\eqref{eq:C.2.14} holds for any derivative operator
\(\nabla_a\).

Finally, we already remarked in Section~C.1 above that if \(\phi:M\to M\) is a
diffeomorphism, then \((M,g_{ab})\) and \((M,\phi^*g_{ab})\) represent the same
physical spacetime. If we consider a one-parameter family of spacetimes
\((M,g_{ab}(\lambda))\), then \((M,\phi_\lambda^*g_{ab}(\lambda))\) represents the
same physical one-parameter family, where \(\phi_\lambda\) is an arbitrary
one-parameter group of diffeomorphisms. If, as in Sections~4.4 and~7.5, we consider
the first-order perturbation of \(g_{ab}\big\rvert_{\lambda=0}\) obtained by
differentiating \(g_{ab}(\lambda)\) with respect to \(\lambda\) at \(\lambda=0\), we
find that
\(h_{ab}=dg_{ab}/d\lambda\big\rvert_{\lambda=0}\) and
\(h'_{ab}=d(\phi_\lambda^*g_{ab})/d\lambda\big\rvert_{\lambda=0}\) represent the
same physical perturbation. But, it is not difficult to see that
\begin{equation}
  h'_{ab}=h_{ab}-\Lie_vg_{ab},
  \tag{C.2.15}\label{eq:C.2.15}
\end{equation}
where \(v^a\) is the vector field which generates \(\phi_\lambda\) and
\(g_{ab}=g_{ab}(\lambda=0)\). Thus, the gauge freedom in perturbations, \(h_{ab}\),
is given by \(\Lie_vg_{ab}\), where \(v^a\) is an arbitrary vector field.
Furthermore, by Eq.~\eqref{eq:C.2.14} we have
\begin{equation}
\begin{aligned}
  \Lie_vg_{ab}
  &=v^c\nabla_cg_{ab}+g_{cb}\nabla_av^c+g_{ac}\nabla_bv^c\\
  &=\nabla_av_b+\nabla_bv_a,
\end{aligned}
  \tag{C.2.16}\label{eq:C.2.16}
\end{equation}
where the second line of Eq.~\eqref{eq:C.2.16} holds when \(\nabla_a\) is the
derivative operator associated with \(g_{ab}\). Thus, the gauge transformations of
linearized general relativity about a solution \(g_{ab}\) are
\begin{equation}
  h_{ab}\longrightarrow h'_{ab}
  =h_{ab}-\nabla_av_b-\nabla_bv_a.
  \tag{C.2.17}\label{eq:C.2.17}
\end{equation}
This is closely analogous to the gauge freedom
\(A_a\longrightarrow A'_a=A_a-\nabla_a\chi\) of electromagnetism.
